%-------------------------
% Cover Letter Template (Block Letter, ATS-Friendly)
% Single-Column, No Section Rules, Real Paragraph Breaks
%-------------------------

\documentclass[letterpaper,11pt]{article}

\usepackage{xcolor}
\usepackage{verbatim}
\usepackage[hidelinks]{hyperref}
\usepackage[margin=1in]{geometry}

% Searchable, faithful text extraction
\usepackage[T1]{fontenc}
\usepackage{lmodern}
\usepackage{mathptmx} % Times body face (core MiKTeX fonts only; small-caps fall back to regular)
\input{glyphtounicode}
\pdfgentounicode=1

\pagestyle{empty}

% URL style
\urlstyle{same}

\raggedbottom
\setlength{\tabcolsep}{0in}
\setlength{\parskip}{6pt}
\setlength{\parindent}{0pt}

% No hyphenation: keeps every word intact for text extraction.
\hyphenpenalty=10000
\exhyphenpenalty=10000
\tolerance=2000
\emergencystretch=1em

%-------------------------------------------
%%%%%%  LETTER STARTS HERE  %%%%%%%%%%%%%%%%%

\begin{document}


\begin{center}
{\Large \bfseries Darshan Maniar} \\ \vspace{2pt}
{\small Heilbronn, Germany | drshnmaniar@gmail.com | linkedin.com/in/drshnmaniar | github.com/drshnmaniar}
\end{center}
\begin{flushright}
{\small \today}
\end{flushright}


{\small \textbf{Subject:} Application for Working Student (f/m/d) eMobility Software (Job ID 519321)}


{\small Dear eMobility Software Team,}


{\small I am applying for the Working Student position in eMobility Software (Job ID 519321) because the role sits exactly at the intersection I work in every day: building dashboards and using AI tools to understand how software and assets actually perform in production. The chance to support project coordination while contributing ideas to innovative eMobility solutions is what drew me in — I enjoy the mix of analytical work and proactive stakeholder communication your posting describes.}


{\small In my current role as Software Engineer III at Reveation Labs, I architected a real-time health monitoring and dashboarding system with React, .NET microservices, and Azure Workbooks across more than 50 global applications, cutting mean time to failure detection from 8–10 hours to 10–15 minutes and shifting support from reactive complaints to proactive fixes. Separately, my first-semester M.Sc. research project at Hochschule Heilbronn was a 29-day structured study on AI tool productivity (36 responses, analyzed with Python-based quantitative and thematic methods), which found 78\% of researchers reported gains while surfacing reliability gaps that eroded trust — hands-on experience evaluating AI tools critically rather than taking their output at face value. On the data side, I have production-grade SQL experience from owning backend services and complex stored procedures sustaining 50K peak orders, plus ETL pipeline work feeding a long-term repository across 150–180 partner clinics. I also led PI planning sessions and mentored junior developers, which built the organized, communicative coordination habits needed to support ongoing projects.}


{\small As an M.Sc. Software Engineering \& Management student at Hochschule Heilbronn (2025–early 2027), I meet the enrollment requirement and am available up to 20 hours per week during the semester for hybrid work around Erlangen. I bring fluent English, German at A2 and actively improving, and a genuine interest in electric mobility as a domain where software quality directly shapes sustainability outcomes. I would welcome the chance to discuss how my dashboard, SQL/Python, and AI-tool experience can support your team's software performance analysis — thank you for your consideration.}



\vspace{6pt}
{\small Best regards,\\}
{\small Darshan Maniar}



\end{document}
